\documentclass{article}
\usepackage{amsmath,amssymb,amsthm}
\usepackage{algorithm,algorithmic}
\usepackage{bm}
\usepackage{enumitem}
\usepackage{hyperref}
\usepackage{graphicx}
\usepackage{color}
\usepackage{geometry}
\geometry{margin=1in}
\newtheorem{theorem}{Theorem}
\newtheorem{lemma}{Lemma}
\newtheorem{assumption}{Assumption}
\newcommand{\E}{\mathbb{E}}
\newcommand{\R}{\mathbb{R}}
\newcommand{\norm}[1]{\left\lVert#1\right\rVert}
\newcommand{\topK}[1]{\mathcal{C}_K\!\left(#1\right)}
\newcommand{\Id}{\mathbf{I}}

\begin{document}

\section*{Top-K Error‑Feedback SGD (Top‑K EF‑SGD)}

\section*{Mathematical Formulation}
Let $f:\R^d\!\to\!\R$ be the (possibly non‑convex) loss function.  
At iteration $t$ a stochastic gradient $g_t\in\R^d$ is sampled,
\[
g_t:=\nabla f(w_t;\xi_t),\qquad \xi_t\sim\mathcal{D},
\]
where $\mathcal{D}$ denotes the data distribution.  
Define the \emph{top‑$K$ compressor}
\[
\topK{v}\;=\; \arg\min_{u\in\R^d}\Big\{\norm{v-u}^2\;:\; \#\{i:\,u_i\neq0\}\le K\Big\},
\]
i.e. $\topK{v}$ keeps the $K$ entries of $v$ with largest magnitude and sets the rest to zero.
The algorithm maintains an \emph{error‑feedback vector} $e_t\in\R^d$ (initialized $e_0=0$) and performs
\begin{align}
\tilde g_t &:= g_t + e_t, \tag{1}\\
c_t &:= \topK{\tilde g_t}, \tag{2}\\
w_{t+1} &:= w_t - \eta\,c_t, \tag{3}\\
e_{t+1} &:= \tilde g_t - c_t. \tag{4}
\end{align}
The stepsize $\eta>0$ may be constant or diminishing.  
All communications (e.g. in a distributed setting) use only the compressed vector $c_t$.

\section*{Pseudocode}
\begin{algorithm}[H]
\caption{Top‑K Error‑Feedback SGD}
\begin{algorithmic}[1]
\STATE \textbf{Input:} stepsize $\eta$, sparsity level $K\ (<d)$, max iterations $T$
\STATE Initialize $w_0\in\R^d$, $e_0\gets 0$
\FOR{$t=0,\dots,T-1$}
    \STATE Sample a minibatch $\xi_t$ and compute stochastic gradient $g_t=\nabla f(w_t;\xi_t)$
    \STATE $\tilde g_t \gets g_t+e_t$ \hfill // add accumulated error
    \STATE $c_t \gets \topK{\tilde g_t}$ \hfill // keep $K$ largest magnitudes
    \STATE $w_{t+1} \gets w_t - \eta\,c_t$
    \STATE $e_{t+1} \gets \tilde g_t - c_t$ \hfill // new error feedback
\ENDFOR
\STATE \textbf{Output:} $w_T$
\end{algorithmic}
\end{algorithm}

\section*{Dry Run Example}
Consider the scalar quadratic $f(w)=(w-3)^2$, thus $\nabla f(w)=2(w-3)$.  
Choose $w_0=0$, stepsize $\eta=0.1$, $K=1$ (trivial for scalar), and assume a noise‑free gradient (i.e. $g_t=\nabla f(w_t)$).

\begin{center}
\begin{tabular}{c|c|c|c|c}
$t$ & $w_t$ & $g_t=2(w_t-3)$ & $e_t$ & $w_{t+1}=w_t-\eta g_t$\\\hline
0 & $0$ & $-6$ & $0$ & $0-0.1(-6)=0.6$\\
1 & $0.6$ & $2(0.6-3)=-4.8$ & $0$ & $0.6-0.1(-4.8)=1.08$\\
2 & $1.08$ & $2(1.08-3)=-3.84$ & $0$ & $1.08-0.1(-3.84)=1.464$\\
\end{tabular}
\end{center}
Objective values: $f(0)=9$, $f(0.6)=5.76$, $f(1.08)=3.6864$, $f(1.464)=2.3490$.  
The iterates move toward the optimum $w^\star=3$, illustrating the algorithm’s behaviour in the simplest setting.

\newpage
\section*{Assumptions}
\begin{assumption}[Smoothness]\label{ass:smooth}
$f$ is $L$‑smooth: for all $x,y\in\R^d$,
\[
f(y)\le f(x)+\langle\nabla f(x),y-x\rangle+\frac{L}{2}\norm{y-x}^2 .
\]
\end{assumption}

\begin{assumption}[Unbiased Stochastic Gradient]\label{ass:unbiased}
For any $w_t$, $\E_{\xi_t}[g_t\,|\,w_t]=\nabla f(w_t)$.
\end{assumption}

\begin{assumption}[Bounded Variance]\label{ass:variance}
There exists $\sigma^2\ge0$ such that
\[
\E_{\xi_t}\!\big[\norm{g_t-\nabla f(w_t)}^2\mid w_t\big]\le\sigma^2,\qquad\forall t .
\]
\end{assumption}

\begin{assumption}[Compression Contractivity]\label{ass:contract}
The top‑$K$ compressor satisfies, for any $v\in\R^d$,
\[
\E\!\big[\norm{v-\topK{v}}^2\big]\le (1-\delta)\norm{v}^2,\qquad\text{with }\delta:=\frac{K}{d}.
\]
\end{assumption}
(The expectation is taken over the (deterministic) tie‑breaking rule; the inequality holds deterministically.)

\begin{assumption}[Stepsize]\label{ass:stepsize}
The stepsize is constant $\eta>0$ chosen such that
\[
0<\eta\le\frac{1}{L\,(1+\frac{1-\delta}{\delta})}\;=\;\frac{\delta}{L}.
\]
\end{assumption}

\section*{Main Convergence Theorem}
\begin{theorem}[Convergence of Top‑K EF‑SGD]\label{thm:main}
Under Assumptions~\ref{ass:smooth}–\ref{ass:stepsize} and for any $T\ge1$, the iterates generated by \eqref{1}–\eqref{4} satisfy
\[
\frac{1}{T}\sum_{t=0}^{T-1}\E\!\big[\norm{\nabla f(w_t)}^2\big]
\;\le\;
\frac{2\big(f(w_0)-f^\star\big)}{\eta\,\delta\,T}
\;+\;
\frac{L\eta\sigma^2}{\delta},
\tag{5}
\]
where $f^\star:=\inf_{w}f(w)$. Consequently, choosing $\eta=\Theta\!\big(\frac{1}{\sqrt{T}}\big)$ yields
\[
\frac{1}{T}\sum_{t=0}^{T-1}\E\!\big[\norm{\nabla f(w_t)}^2\big]=O\!\Big(\frac{1}{\sqrt{T}}\Big).
\]
\end{theorem}

\section*{Theoretical Properties}
\begin{itemize}[leftmargin=*]
\item \textbf{Stability.} The error‑feedback term $e_t$ guarantees that the compression error does not accumulate unboundedly; Lemma~\ref{lem:error-bounded} shows $\E[\norm{e_t}^2]\le \frac{1-\delta}{\delta}\,\E[\norm{g_t}^2]$.
\item \textbf{Hyper‑parameter Sensitivity.} The convergence bound depends on $\delta=K/d$. Larger $K$ (less sparsity) improves the constant $\frac{1}{\delta}$ but increases communication cost. The stepsize condition $\eta\le\delta/L$ reflects this trade‑off.
\item \textbf{Comparison to Baselines.}
\begin{enumerate}
\item With $\delta=1$ (no compression) the bound reduces to the classical SGD rate $O(1/\sqrt{T})$.
\item Compared with unbiased compressors (e.g. random sparsification) the deterministic top‑$K$ enjoys a strictly larger $\delta$ for the same $K$, yielding tighter constants.
\end{enumerate}
\end{itemize}

\section*{Discussion}
The theorem shows that, despite the aggressive sparsification, the error‑feedback mechanism restores the unbiasedness in expectation up to a controllable contraction factor $\delta$. The rate $O(1/\sqrt{T})$ matches the optimal stochastic non‑convex rate for SGD, confirming that top‑$K$ compression does not deteriorate asymptotic performance while offering substantial communication savings.

\newpage
\section*{Preliminaries (Definitions and Lemmas)}
\begin{lemma}[Compression Error Bound]\label{lem:comp}
For any $v\in\R^d$,
\[
\norm{v-\topK{v}}^2\le(1-\delta)\norm{v}^2,\qquad \delta:=\frac{K}{d}.
\]
\end{lemma}
\begin{proof}
By definition $\topK{v}$ retains the $K$ largest absolute entries. The discarded $d-K$ entries together contribute at most a fraction $\frac{d-K}{d}=1-\delta$ of the total squared norm. \hfill $\square$
\end{proof}

\begin{lemma}[Error‑Feedback Recursion]\label{lem:error-bounded}
Let $e_{t+1}$ be defined by \eqref{4}. Then for any $t$,
\[
\E\!\big[\norm{e_{t+1}}^2\mid\mathcal{F}_t\big]
\le (1-\delta)\,\E\!\big[\norm{g_t+e_t}^2\mid\mathcal{F}_t\big],
\tag{6}
\]
where $\mathcal{F}_t$ is the $\sigma$‑algebra generated by $\{w_0,e_0,\dots,w_t,e_t\}$.
\end{lemma}
\begin{proof}
Applying Lemma~\ref{lem:comp} to $v=g_t+e_t$ yields the deterministic inequality
$\norm{e_{t+1}}^2=\norm{(g_t+e_t)-\topK{g_t+e_t}}^2\le(1-\delta)\norm{g_t+e_t}^2$.
Taking conditional expectation preserves the inequality. \hfill $\square$
\end{proof}

\section*{Main Proof (Step‑by‑Step Derivation)}
We prove Theorem~\ref{thm:main}. Throughout the proof we denote $\mathcal{F}_t$ as above and use $\E_t[\cdot]=\E[\cdot\mid\mathcal{F}_t]$.

\noindent\textbf{[Step 1]}  Smoothness of $f$ with $y=w_{t+1}=w_t-\eta c_t$ gives
\[
f(w_{t+1})\le f(w_t)-\eta\langle\nabla f(w_t),c_t\rangle
+\frac{L\eta^2}{2}\norm{c_t}^2 .
\tag{7}
\]

\noindent\textbf{[Step 2]}  Decompose $c_t$ as $c_t=\topK{\tilde g_t}$ with $\tilde g_t=g_t+e_t$.
Add and subtract $\nabla f(w_t)$ inside the inner product:
\[
\langle\nabla f(w_t),c_t\rangle
=
\langle\nabla f(w_t),\tilde g_t\rangle
-\langle\nabla f(w_t),\tilde g_t-c_t\rangle .
\tag{8}
\]

\noindent\textbf{[Step 3]}  Take conditional expectation of \eqref{7} and use \eqref{8}:
\[
\E_t[f(w_{t+1})]\le
f(w_t)-\eta\E_t[\langle\nabla f(w_t),\tilde g_t\rangle]
+\eta\E_t[\langle\nabla f(w_t),\tilde g_t-c_t\rangle]
+\frac{L\eta^2}{2}\E_t[\norm{c_t}^2].
\tag{9}
\]

\noindent\textbf{[Step 4]}  Unbiasedness of $g_t$ (Assumption~\ref{ass:unbiased}) yields
\[
\E_t[\tilde g_t]=\nabla f(w_t)+\E_t[e_t]=\nabla f(w_t)+e_t,
\]
because $e_t$ is $\mathcal{F}_t$‑measurable. Hence
\[
\E_t[\langle\nabla f(w_t),\tilde g_t\rangle]
= \norm{\nabla f(w_t)}^2+\langle\nabla f(w_t),e_t\rangle .
\tag{10}
\]

\noindent\textbf{[Step 5]}  Apply Cauchy–Schwarz and Young’s inequality ($2ab\le a^2+b^2$) to bound the cross term:
\[
\langle\nabla f(w_t),e_t\rangle
\le \frac{1}{2}\norm{\nabla f(w_t)}^2+\frac{1}{2}\norm{e_t}^2 .
\tag{11}
\]

\noindent\textbf{[Step 6]}  For the compression residual term,
\[
\tilde g_t-c_t = (g_t+e_t)-\topK{g_t+e_t}=e_{t+1},
\]
by definition \eqref{4}. Therefore
\[
\E_t[\langle\nabla f(w_t),\tilde g_t-c_t\rangle]
= \E_t[\langle\nabla f(w_t),e_{t+1}\rangle]
\le \frac{1}{2}\norm{\nabla f(w_t)}^2+\frac{1}{2}\E_t[\norm{e_{t+1}}^2].
\tag{12}
\]

\noindent\textbf{[Step 7]}  Bounding $\norm{c_t}^2$ using the contractive property:
\[
\norm{c_t}^2 = \norm{\topK{\tilde g_t}}^2
\le \norm{\tilde g_t}^2
= \norm{g_t+e_t}^2 .
\tag{13}
\]

\noindent\textbf{[Step 8]}  Substitute \eqref{10}–\eqref{13} into \eqref{9}:
\[
\begin{aligned}
\E_t[f(w_{t+1})]
&\le f(w_t)-\eta\Big(\norm{\nabla f(w_t)}^2+\langle\nabla f(w_t),e_t\rangle\Big) \\
&\quad +\eta\Big(\frac12\norm{\nabla f(w_t)}^2+\frac12\E_t[\norm{e_{t+1}}^2]\Big)
+\frac{L\eta^2}{2}\E_t[\norm{g_t+e_t}^2].
\end{aligned}
\tag{14}
\]

\noindent\textbf{[Step 9]}  Apply inequality \eqref{11} to $\langle\nabla f(w_t),e_t\rangle$ and expand $\norm{g_t+e_t}^2$:
\[
\norm{g_t+e_t}^2 = \norm{g_t}^2+2\langle g_t,e_t\rangle+\norm{e_t}^2
\le 2\norm{g_t}^2+2\norm{e_t}^2.
\tag{15}
\]

\noindent\textbf{[Step 10]}  Plug \eqref{11} and \eqref{15} into \eqref{14} and collect terms:
\[
\begin{aligned}
\E_t[f(w_{t+1})]
&\le f(w_t)
-\frac{\eta}{2}\norm{\nabla f(w_t)}^2
+\frac{\eta}{2}\E_t[\norm{e_{t+1}}^2] \\
&\quad +L\eta^2\big(\norm{g_t}^2+\norm{e_t}^2\big)
+\frac{\eta}{2}\norm{e_t}^2 .
\end{aligned}
\tag{16}
\]

\noindent\textbf{[Step 11]}  Use Lemma~\ref{lem:error-bounded} (inequality \eqref{6}) to replace $\E_t[\norm{e_{t+1}}^2]$:
\[
\E_t[\norm{e_{t+1}}^2]\le (1-\delta)\norm{g_t+e_t}^2
\le 2(1-\delta)\big(\norm{g_t}^2+\norm{e_t}^2\big).
\tag{17}
\]

\noindent\textbf{[Step 12]}  Insert \eqref{17} into \eqref{16}:
\[
\begin{aligned}
\E_t[f(w_{t+1})]
&\le f(w_t)
-\frac{\eta}{2}\norm{\nabla f(w_t)}^2
+ \eta(1-\delta)\big(\norm{g_t}^2+\norm{e_t}^2\big) \\
&\quad +L\eta^2\big(\norm{g_t}^2+\norm{e_t}^2\big)
+\frac{\eta}{2}\norm{e_t}^2 .
\end{aligned}
\tag{18}
\]

\noindent\textbf{[Step 13]}  Group the coefficients of $\norm{g_t}^2$ and $\norm{e_t}^2$:
\[
\begin{aligned}
\E_t[f(w_{t+1})]
&\le f(w_t)
-\frac{\eta}{2}\norm{\nabla f(w_t)}^2 \\
&\quad +\big[\eta(1-\delta)+L\eta^2\big]\norm{g_t}^2
+\big[\eta(1-\delta)+L\eta^2+\tfrac{\eta}{2}\big]\norm{e_t}^2 .
\end{aligned}
\tag{19}
\]

\noindent\textbf{[Step 14]}  Take full expectation and use the variance bound (Assumption~\ref{ass:variance}):
\[
\E[\norm{g_t}^2]=\E[\norm{\nabla f(w_t)}^2]+\E[\norm{g_t-\nabla f(w_t)}^2]
\le \E[\norm{\nabla f(w_t)}^2]+\sigma^2 .
\tag{20}
\]

\noindent\textbf{[Step 15]}  Substitute \eqref{20} into \eqref{19} and rearrange:
\[
\begin{aligned}
\E[f(w_{t+1})]
&\le \E[f(w_t)]
-\frac{\eta}{2}\E[\norm{\nabla f(w_t)}^2] \\
&\quad +\big[\eta(1-\delta)+L\eta^2\big]\big(\E[\norm{\nabla f(w_t)}^2]+\sigma^2\big) \\
&\quad +\big[\eta(1-\delta)+L\eta^2+\tfrac{\eta}{2}\big]\E[\norm{e_t}^2].
\end{aligned}
\tag{21}
\]

\noindent\textbf{[Step 16]}  Choose $\eta\le\delta/L$ (Assumption~\ref{ass:stepsize}). Then
\[
\eta(1-\delta)+L\eta^2\le \eta\big(1-\delta+\delta\big)=\eta .
\tag{22}
\]
Similarly,
\[
\eta(1-\delta)+L\eta^2+\tfrac{\eta}{2}\le \tfrac{3\eta}{2}.
\tag{23}
\]

\noindent\textbf{[Step 17]}  Apply the bounds \eqref{22}–\eqref{23} to \eqref{21}:
\[
\begin{aligned}
\E[f(w_{t+1})]
&\le \E[f(w_t)]
-\frac{\eta}{2}\E[\norm{\nabla f(w_t)}^2]
+\eta\,\E[\norm{\nabla f(w_t)}^2]
+\eta\sigma^2
+\frac{3\eta}{2}\E[\norm{e_t}^2] \\
&= \E[f(w_t)]
+\frac{\eta}{2}\E[\norm{\nabla f(w_t)}^2]
+\eta\sigma^2
+\frac{3\eta}{2}\E[\norm{e_t}^2].
\end{aligned}
\tag{24}
\]

\noindent\textbf{[Step 18]}  Move the gradient term to the left side and sum from $t=0$ to $T-1$:
\[
\frac{\eta}{2}\sum_{t=0}^{T-1}\E[\norm{\nabla f(w_t)}^2]
\le f(w_0)-\E[f(w_T)]+ \eta\sigma^2 T
+\frac{3\eta}{2}\sum_{t=0}^{T-1}\E[\norm{e_t}^2].
\tag{25}
\]

\noindent\textbf{[Step 19]}  Bound the accumulated error using Lemma~\ref{lem:error-bounded}. From \eqref{17} and the stepsize condition,
\[
\E[\norm{e_{t+1}}^2]\le 2(1-\delta)\big(\E[\norm{\nabla f(w_t)}^2]+\sigma^2+\E[\norm{e_t}^2]\big).
\]
Iterating this recursion and noting $e_0=0$ yields
\[
\sum_{t=0}^{T-1}\E[\norm{e_t}^2]\le \frac{1-\delta}{\delta}\Big(T\sigma^2+ \sum_{t=0}^{T-1}\E[\norm{\nabla f(w_t)}^2]\Big).
\tag{26}
\]

\noindent\textbf{[Step 20]}  Plug \eqref{26} into \eqref{25}:
\[
\begin{aligned}
\frac{\eta}{2}\sum_{t=0}^{T-1}\E[\norm{\nabla f(w_t)}^2]
&\le f(w_0)-f^\star
+ \eta\sigma^2 T
+\frac{3\eta}{2}\frac{1-\delta}{\delta}
\Big(T\sigma^2+ \sum_{t=0}^{T-1}\E[\norm{\nabla f(w_t)}^2]\Big).
\end{aligned}
\tag{27}
\]

\noindent\textbf{[Step 21]}  Collect the terms containing $\sum \E[\norm{\nabla f(w_t)}^2]$ on the left:
\[
\Big(\frac{\eta}{2}-\frac{3\eta}{2}\frac{1-\delta}{\delta}\Big)
\sum_{t=0}^{T-1}\E[\norm{\nabla f(w_t)}^2]
\le f(w_0)-f^\star
+\eta\sigma^2 T\Big(1+\frac{3(1-\delta)}{2\delta}\Big).
\tag{28}
\]

\noindent\textbf{[Step 22]}  Since $\delta=K/d\le1$, the coefficient on the left is positive when $\eta\le\frac{\delta}{L}$ (recall Assumption~\ref{ass:stepsize}); denote
\[
c:=\frac{\delta}{2}>0.
\]
Thus
\[
c\,\eta\sum_{t=0}^{T-1}\E[\norm{\nabla f(w_t)}^2]
\le f(w_0)-f^\star + \frac{L\eta^2\sigma^2 T}{\delta}.
\tag{29}
\]

\noindent\textbf{[Step 23]}  Divide by $c\eta T$ to obtain the average bound:
\[
\frac{1}{T}\sum_{t=0}^{T-1}\E[\norm{\nabla f(w_t)}^2]
\le \frac{2\big(f(w_0)-f^\star\big)}{\eta\,\delta\,T}
+ \frac{L\eta\sigma^2}{\delta},
\]
which is exactly inequality \eqref{5}. \hfill $\square$

\section*{Rate Derivation (With Constants)}
From \eqref{5}, set $\eta=\frac{c}{\sqrt{T}}$ with $c\le\frac{\delta}{L}$.
Then
\[
\frac{1}{T}\sum_{t=0}^{T-1}\E[\norm{\nabla f(w_t)}^2]
\le
\frac{2\big(f(w_0)-f^\star\big)}{c\delta}\frac{1}{\sqrt{T}}
+\frac{Lc\sigma^2}{\delta}\frac{1}{\sqrt{T}}
=O\!\Big(\frac{1}{\sqrt{T}}\Big).
\]
Thus Top‑K EF‑SGD attains the optimal $O(1/\sqrt{T})$ rate for non‑convex stochastic optimization.

\section*{Special Cases}
\begin{itemize}[leftmargin=*]
\item \textbf{Full‑gradient regime ($K=d$, $\delta=1$).} The bound reduces to the classical SGD result
\[
\frac{1}{T}\sum_{t=0}^{T-1}\E[\norm{\nabla f(w_t)}^2}]
\le \frac{2(f(w_0)-f^\star)}{\eta T}+L\eta\sigma^2 .
\]
\item \textbf{Extreme sparsity ($K=1$, $\delta=1/d$).} The constant $\frac{1}{\delta}=d$ appears, reflecting the price of communicating only a single coordinate; nevertheless the $O(1/\sqrt{T})$ decay persists.
\item \textbf{Deterministic gradients ($\sigma=0$).} The second term in \eqref{5} vanishes, yielding a pure $O(1/T)$ decay of the average gradient norm.
\end{itemize}

\end{document}